\section*{Capítulo \ref{ch:g6:light-propagation} --- Propagación rectilínea de la luz}

\begin{solution}{exo:g6:light-propagation:1}
En \omterm{def:g2:air-around-us:air}{aire} limpio (o en cualquier material transparente y homogéneo) la
\omterm{def:g1:light-and-shadows:source}{luz} viaja en línea recta. \omterm{def:g6:describing-motion:trajectory}{Rectilínea}: en línea recta;
propagación: el viaje de la \omterm{def:g1:light-and-shadows:source}{luz} de un sitio a otro.
\end{solution}

\begin{solution}{exo:g6:light-propagation:2}
Divergente: rayos que se abren desde un punto (una vela, una
\omterm{def:g3:first-electric-circuit:bulb}{bombilla} desnuda). Paralelo: rayos uno al lado del otro (la \omterm{def:g1:light-and-shadows:source}{luz} del
Sol --- el Sol está tan lejos que sus rayos nos llegan desfilando
juntos). Convergente: rayos que se cierran hacia un punto (\omterm{def:g1:light-and-shadows:source}{luz}
recogida por un instrumento de aumento).
\end{solution}

\begin{solution}{exo:g6:light-propagation:3}
Hace de regla un \omterm{def:g6:light-propagation:ray}{rayo de luz}: el ojo, el canto y el extremo lejano
se juzgan alineados exactamente cuando la \omterm{def:g1:light-and-shadows:source}{luz} del extremo lejano
llega al ojo rozando el canto entero --- la
\omterm{prop:g6:light-propagation:law}{propagación rectilínea} es el juramento.
\end{solution}

\begin{solution}{exo:g6:light-propagation:4}
Los puntos de aterrizaje de los rayos que rozan justo los bordes del
obstáculo. Un ojo colocado dentro del cono de \omterm{def:g1:light-and-shadows:shadow}{sombra} no puede ver la
lámpara --- no le llega ningún rayo recto salido de ella.
\end{solution}

\begin{solution}{exo:g6:light-propagation:5}
A mitad de camino: la \omterm{def:g1:light-and-shadows:shadow}{sombra} es el doble de la pelota ---
\qty{12}{cm}. Cerca de la pantalla: la \omterm{def:g1:light-and-shadows:shadow}{sombra} encoge hacia los
\qty{6}{cm} verdaderos de la pelota.
\end{solution}

\begin{solution}{exo:g6:light-propagation:6}
En la parte de abajo de la pantalla: una línea recta que viene de un
punto alto y pasa por la puerta baja tiene que seguir bajando.
\end{solution}

\begin{solution}{exo:g6:light-propagation:7}
Del revés y con la izquierda y la derecha intercambiadas: en la
pantalla la mano se agita abajo en vez de arriba, y del lado
contrario. Las dos inversiones vienen del cruce de los hilos en la
puerta.
\end{solution}

\begin{solution}{exo:g6:light-propagation:8}
Astrónomos: \omterm{def:g1:five-senses:observation}{observar} los eclipses sin peligro en la pantalla en vez de
en el cielo. Pintores: calcar la escena proyectada para dominar la
perspectiva. (Parientes modernos: todas las cámaras y el propio ojo.)
\end{solution}

\begin{solution}{exo:g6:light-propagation:9}
La \omterm{def:g1:light-and-shadows:shadow}{sombra} encoge hacia el tamaño verdadero de la moneda: con la
lámpara más lejos, los rayos rasantes salen de ella con menos
inclinación --- un abanico más estrecho --- y aterrizan más juntos. La
construcción lo enseña: líneas más tendidas, apertura más ceñida.
\end{solution}

\begin{solution}{exo:g6:light-propagation:10}
El Sol es compacto y lejano: un único abanico ceñido de rayos casi
paralelos, una única frontera nítida --- \omterm{def:g1:light-and-shadows:shadow}{sombras} bien marcadas. Una
capa de nubes grises es una lámpara del tamaño del cielo entero: la
\omterm{def:g1:light-and-shadows:source}{luz} llega de todas las direcciones a la vez, a cada punto del suelo
le llega algo de ella y las \omterm{def:g1:light-and-shadows:shadow}{sombras} se desdibujan casi hasta la
nada.
\end{solution}

\begin{solution}{exo:g6:light-propagation:11}
Más brillante: una puerta ancha deja entrar más hilos de cada punto de
la escena. Emborronada: a través de una puerta ancha, los hilos de
cada punto de la escena llegan a toda una \emph{mancha} de pantalla en
vez de a un solo punto --- e incontables manchas superpuestas
emborronan la \omterm{prop:g5:mirrors-reflection:image}{imagen}. La nitidez del agujerito \emph{es} su pequeñez.
\end{solution}

\begin{solution}{exo:g6:light-propagation:12}
Cada hueco entre las hojas es un agujerito natural; cada mancha
brillante del suelo es la \omterm{prop:g5:mirrors-reflection:image}{imagen} del Sol que forma ese agujerito
--- redonda un día corriente porque el Sol es redondo. Cuando el
eclipse convierte al Sol en un creciente, cada hueco entre hojas
proyecta fielmente el creciente: cientos de camaritas, un solo
protagonista celeste.
\end{solution}

\begin{solution}{pb:g6:light-propagation:1}
\textbf{1.} Pegada a la pantalla, a los rayos rasantes no les queda
distancia para abrirse: la \omterm{def:g1:light-and-shadows:shadow}{sombra} se ciñe al contorno de la propia
marioneta.
\textbf{2.} A mitad de camino se dobla: \qty{40}{cm}.
\textbf{3.} Cuatro veces $20$ son exactamente \qty{80}{cm} --- cuadra;
la marioneta está a un cuarto de la distancia lámpara-pantalla.
\textbf{4.} Un globo luminoso ancho son muchas lámparas a la vez:
cada parte echa su propia \omterm{def:g1:light-and-shadows:shadow}{sombra} ligeramente desplazada, y los
bordes del dragón salen emborronados y grises. El teatro nítido
necesita una fuente pequeña, desnuda y casi puntual.
\textbf{5.} Hacia la lámpara: el abanico de rayos rasantes se abre y
el lobo crece; el brazo se queda abajo, fuera del abanico. Según la
marioneta se acerca a la lámpara, pequeña pero real, los bordes se
ablandan un poco más --- de cerca, el ancho de la lámpara importa
más.
\textbf{6.} Las líneas rectas que van de la lámpara a cada marioneta
tienen que llegar a regiones separadas de la pantalla --- las
marionetas están lo bastante separadas de lado como para que ninguna
quede dentro del cono de \omterm{def:g1:light-and-shadows:shadow}{sombra} de la otra.
\textbf{7.} La de la mano: al estar más cerca de la lámpara, a un
tercio del camino, sus rayos rasantes se abren tres veces, mientras
que los del dragón, a mitad de escena, solo se doblan --- cerca de la
lámpara, hasta las manos pequeñas hacen de gigantes.
\textbf{8.} Llevar al villano pegado a la pantalla, donde encoge
hasta su tamaño normal (o sacarlo del \omterm{def:g6:light-propagation:ray}{haz} por completo); o colar
un segundo obstáculo cerca de la lámpara --- su \omterm{def:g1:light-and-shadows:shadow}{sombra} gigante se
traga la del villano, y el villano ``desaparece'' dentro de otra
oscuridad.
\textbf{9.} Abajo en la sábana: desde la lámpara alta, el hilo recto
que pasa por el agujero sigue bajando --- el agujerito invierte.
\textbf{10.} De derecha a izquierda: la puerta intercambia la
izquierda y la derecha junto con el arriba y el abajo.
\textbf{11.} Más brillante, pero desesperadamente borrosa: cada punto
de la calle pinta ahora una mancha del tamaño de una moneda, las
manchas se superponen y la \omterm{prop:g5:mirrors-reflection:image}{imagen} se deshace --- el final muere junto
con la nitidez.
\textbf{12.} ``El programa entero de esta noche --- todas las
\omterm{def:g1:light-and-shadows:shadow}{sombras}, todos los gigantes y una calle del revés --- lo ha
interpretado una sola ley: la \omterm{def:g1:light-and-shadows:source}{luz} viaja en línea recta''.
\end{solution}
